% !TEX TS-program = pdflatex
% !TEX encoding = UTF-8 Unicode
\documentclass[a4paper,twoside,11pt]{article}
\usepackage[polish]{babel}
\usepackage[MeX]{polski}
\usepackage[utf8]{inputenc}
\usepackage{indentfirst,graphicx,amsmath}
\usepackage[left=2cm,top=2cm,right=2cm,nohead]{geometry}
\usepackage{hyperref}
\usepackage{tabularx}
\usepackage{color}
\usepackage{float}
\frenchspacing
\pagestyle{plain}

\makeatletter
\renewcommand\@seccntformat[1]{\csname the#1\endcsname.\quad}
\renewcommand\numberline[1]{#1.\hskip0.7em}
\makeatother

\newenvironment{mylisting}
{\begin{list}{}{\setlength{\leftmargin}{1em}}\item\scriptsize\bfseries}
{\end{list}}

\newenvironment{mytinylisting}
{\begin{list}{}{\setlength{\leftmargin}{1em}}\item\tiny\bfseries}
{\end{list}}

\usepackage{graphicx}
\usepackage{color}

\author{\\ ~ \\ Tomasz Gański \\
Tomasz Gieniusz \\
Bartosz Jankowski\\
Grzegorz Marcinkowski\\
Łukasz Odzioba\\
Jacek Weremko}
\title{\LARGE Wykrywacz zachcianek gastronomicznych}

\begin{document}

\makeatletter

\renewcommand{\maketitle}{
	\begin{titlepage}
		\begin{center}
			\small Politechnika Gdańska \\ Wydział Elektroniki, Telekomunikacji i Informatyki
		\end{center}
		\vspace{3cm}
		\noindent \rule{\linewidth}{0.4mm}
		\begin{center}
			\textsc{
				\huge{Architektura systemu } \\
				\vspace{0.3ex}
				\LARGE{Systemy z Bazą Wiedzy}
			}
			\\
			\vspace{1em}
			\Large{\@title}
		\end{center}
		\rule{\linewidth}{0.4mm}
		\vspace{1em}
		\begin{center}
			\large Autorzy: \@author
		\end{center}
		\vspace{7em}
		\ifpdf
			\begin{tabularx}{\textwidth}{ X X }
				\makebox[0.5\textwidth][l]{
					\hspace{1cm}
					\includegraphics[width=4cm]{logo_pg}
				}
				&
				\makebox[0.5\textwidth][r]{
					\includegraphics[width=4cm]{logo_eti}
					\hspace{1cm}
				}
			\end{tabularx}
		\fi
		% Duży odstęp, który powoduje, że stopka z datą pojawia się na dole. Dokładniej wypycha ją
		% poza stronę, ale latex na to nie pozwala i umieszcza ją na końcu strony.
		\vspace*{\stretch{6}}
		\begin{center}
			Gdańsk, \@date \ r.
		\end{center}
	\end{titlepage}
}

\makeatother

\maketitle{}


\tableofcontents

\newpage

\section{Cel projektu}
Celem projektu jest stworzenie systemu doradczego, ułatwiającego życie każdego człowieka. Każdy z nas bywa w sytuacjach, gdy jest głodny lecz niezdecydowany co do konkretnej potrawy. \\ 

System na podstawie wielu parametrów, takich jak wiek użytkownika, aktualna pora dnia, roku, pogoda, itp. będzie w stanie ułatwić użytkownikowi podjęcie decyzji sugerując mu typowe potrawy, które mogą pobudzić jego kubki smakowe.

\section{Grupa docelowa}
Grupą docelową są przede wszystkim amerykańskie nastolatki, które z nudów ściągają każdą aplikację dostępną w internecie, a także ludzie regularnie podjadający rozmaite smakołyki i znudzeni niezmienną zawartością szafki ze słodyczami. \\

Dodatkowo aplikacja może zainteresować ludzi lubiących wyzwania, a także przydać się na różnego rodzaju imprezach dostarczając rozrywki.

\section{Wymagania}
\subsection{Wejście do systemu}
System automatycznego doradcy gastronomicznego powinien umożliwiać generowanie podpowiedzi, dotyczących preferencji smakowych użytkownika na podstawie następujących parametrów:
\begin{itemize}
\item pora roku
\item pora dnia
\item pogoda
\item wiek użytkownika
\item płeć użytkownika
\item lokacja użytkownika
\end{itemize}

Opcjonalnie użytkownik może zapewnić dodatkowe dane wejściowe do systemu poruszając telefonem, które wydobędziemy za pomocą akcelerometru. \\

System może także brać pod uwagę opinie użytkownika dotyczącą poprzedniej propozycji. Przykładowo po wykryciu ciągu komentarzy \emph{zbyt tuczące} system domyśli się, że użytkownik jest na diecie i rzadziej będzie proponował tłuste jedzenie. \\

Dodatkową informacją wejściową dla systemu może być informacja co dziś zjadł, tak aby system mógł obliczyć zapotrzebowanie energetyczne użytkownika lub nie proponować użytkownikowi potrawy którą dziś już jadł.

\newpage
\subsection{Wynik}
Informacje wyświetlane po ustaleniu gastronomicznej zachcianki użytkownika:
\begin{itemize}
\item nazwa obiektu pożądania
\item obrazek wzmacniający głód i utwierdzający w słuszności predykcji
\item opis produktu
\end{itemize}

Dzięki wykorzystaniu informacji o położeniu i użyciu wyszukiwarki Google system może zaproponować użytkownikowi
najbliższą pizzerię, gdy ten ma ochotę na pizze, lub w podobny sposób rozwiązać jego problem o którym jeszcze nie wie. Narzędzie to będzie pomocne zwłaszcza po wykryciu, że użytkownik jest w \emph{obcym} mieście
i na pewno doceni pomoc systemu.
\subsection{Dodatkowe}
Dodatkowo aplikacja powinna umożliwiać dzielenie się swoimi zachciankami na Facebooku oraz podglądanie na co mają ochotę nasi znajomi.

Jeśli kilku znajomych ma ochotę na pizze to warto zaproponować im wspólne wyjście do pizzerii, a dzięki współpracy z firmą Groupon możemy uzyskać środki na rozwój systemu.
\section{Technologia}
Projekt zostanie wykonany na platformę mobilną Android (zgodność z wersjami Donut i nowszymi). Wykorzystane zostanie natywne API systemu, a także system AdMob w celu wyświetlania reklam po umieszczeniu aplikacji w Android Market. \\

Aplikacja do działania będzie wymagała dostępu do internetu oraz aktualnej lokacji użytkownika. \\

System uczący będzie zaimplementowany w formie sieci neuronowej której współczynniki będa zapisywane do pamięci trwałej, tak aby zachować jej stan pomiędzy uruchomieniami aplikacji. \\


Interfejs użytkownika powinien być wyraźny i czytelny nawet dla użytkownika znajdującego sie w stanie upojenia alkoholowego.

\section{Zadania do wykonania}
\begin{itemize}
\item analiza zagadnienia sieci neuronowych
\item stworzenie sieci neuronowej
\item przygotowanie zestawu uczacego do testow sieci neuronowej
\item stworzenie bazy produktow
\item opracowanie interfejsu graficznego
\item implementacja systemu powiadomien na facebooku
\item implementacja systemu grupowania zachciianek przyjaciol
\item analiza android API
\item projekt systemu
\item przygotowanie prezentacji produktu
\end{itemize}

\section{Podział pracy}
\begin{center}
\begin{tabular}{|l|p{11cm}|l|}
  \hline 
  Osoba & Realizowane zadania & Czas [h]\\
  \hline
  Łukasz Odzioba & Opracowanie koncepcji & 1 \\
 \hline
  Tomasz Gieniusz & Opracowanie koncepcji & 1 \\
 \hline
  Tomasz Gański & Opracowanie koncepcji & 1 \\
 \hline 
  Bartosz Jankowski & Opracowanie koncepcji & 1 \\
 \hline
  Grzegorz Marcinkowski &  & 1/6 \\
 \hline
  Jacek Weremko &  & 0 \\
 \hline
\end{tabular} 
\end{center}
\end{document}
